La asignatura \textit{Fundamentos de Interpretación y Compilación de Lenguajes de Programación} (FLP) del programa de Ingeniería de Sistemas de la Universidad del Valle sede Tuluá presenta un nivel de abstracción que dificulta su aprendizaje, agravado por la fricción sintáctica del entorno de trabajo basado en Racket y la librería \textit{Essentials of Programming Languages} (EOPL). Un diagnóstico realizado con 36 estudiantes y 5 docentes del programa confirmó esta situación: el 77.8\,\% de los estudiantes reportó dificultades para comprender conceptos como ambientes de ejecución y estado, apenas el 19.4\,\% se sintió preparado para aplicar lo aprendido en otros contextos, y el 94.5\,\% consideró que recursos interactivos adicionales habrían mejorado su aprendizaje.

Con el propósito de atender esta problemática, este trabajo de grado desarrolla FLP Viewer, un prototipo web que apoya la comprensión de los contenidos de la asignatura mediante la visualización de estructuras internas y el uso de lenguajes orientados por sintaxis. El desarrollo siguió un proceso estructurado en cuatro etapas: la caracterización de las dificultades conceptuales a partir del diagnóstico con estudiantes, docentes y la revisión del microcurrículo; la construcción de un marco teórico sobre interpretación de lenguajes, teorías del aprendizaje y usabilidad; el diseño de la arquitectura, las tecnologías y las estrategias de visualización del sistema; y la implementación de un prototipo funcional que permite al estudiante definir gramáticas en notación BNF, visualizar de forma interactiva el árbol de sintaxis abstracta generado y observar la evolución del ambiente de ejecución durante la evaluación paso a paso de un programa.

El prototipo se evaluó con 12 estudiantes de la asignatura y el docente responsable mediante el cuestionario System Usability Scale (SUS) y un instrumento de utilidad pedagógica de elaboración propia. Los resultados obtenidos, un puntaje SUS promedio de 82.7 sobre 100 (categoría ``Excelente'') y una recomendación del prototipo por parte del 91.7\,\% de los participantes, indican que la herramienta es usable y percibida como útil para el aprendizaje de los conceptos de mayor dificultad de la asignatura, en particular la visualización simultánea de la gramática, el AST y el ambiente de ejecución.

\vspace{1em}
\noindent\textbf{Palabras clave:} interpretación de lenguajes de programación, visualización de programas, ambiente de ejecución, árbol de sintaxis abstracta, herramientas educativas interactivas, EOPL, Racket.
